% Приложения (графический материал и листинги программ свыше одной страницы)
% Обозначаются заглавными буквами русского алфавита, за исключением
% букв Ё, З, Й, О, Ч, Ъ, Ы, Ь.

\clearpage
\begin{center}
\textbf{ПРИЛОЖЕНИЕ~А}\\[6pt]
Графический материал
\end{center}
\addcontentsline{toc}{section}{ПРИЛОЖЕНИЕ А. Графический материал}

\vspace{12pt}
% TODO: вставить графический материал по списку (8--12 листов А1 / А4 с рамкой):
% 1) Схема архитектуры приложения
% 2) Диаграмма вариантов использования
% 3) ER-диаграмма базы данных
% 4) Диаграмма последовательности алгоритма отбора абитуриентов
% 5) Блок-схема алгоритма полного пересчёта результатов отбора
% 6) Диаграмма ролевой модели доступа
% 7) Диаграмма архитектуры системы аудита и логического удаления
% 8) Главное окно приложения
% 9) Окно результатов конкурсного списка
% 10) Окно журнала аудита

\noindent\textit{Графический материал оформляется в виде плакатов формата~А4 с рамкой и основной надписью (угловым штампом) в соответствии с приложениями~7 и~8 методички и помещается после списка использованной литературы. Каждый плакат имеет название, обозначение вида \mbox{ДП--1070122218--2026--XX} и сквозную нумерацию.}

\clearpage
\begin{center}
\textbf{ПРИЛОЖЕНИЕ~Б}\\[6pt]
Листинги ключевых модулей программы
\end{center}
\addcontentsline{toc}{section}{ПРИЛОЖЕНИЕ Б. Листинги ключевых модулей}

\vspace{12pt}
\subsection*{Б.1 Сервис JWT-аутентификации (\code{JwtService.cs})}

% TODO: вставить листинг JwtService.cs

\subsection*{Б.2 Ядро алгоритма пересчёта результатов отбора (\code{FullRecalculationCore.cs})}

% TODO: вставить листинг FullRecalculationCore.cs

\subsection*{Б.3 Компаратор кандидатов (\code{CategoryComparer.cs})}

% TODO: вставить листинг CategoryComparer.cs

\subsection*{Б.4 Сервис журналирования (\code{AuditLogger.cs})}

% TODO: вставить листинг AuditLogger.cs

\subsection*{Б.5 Интерцептор Axios на клиенте (\code{api/index.js})}

% TODO: вставить листинг api/index.js
